\documentclass[11pt]{article}

\usepackage[a4paper,margin=23mm]{geometry}
\usepackage[T1]{fontenc}
\usepackage{lmodern}
\usepackage{microtype}
\usepackage{amsmath,amssymb,amsthm,mathtools}
\usepackage{booktabs}
\usepackage{array}
\usepackage[compact]{titlesec}
\usepackage[hidelinks]{hyperref}

\setlength{\abovedisplayskip}{7pt plus 2pt minus 2pt}
\setlength{\belowdisplayskip}{7pt plus 2pt minus 2pt}
\setlength{\abovedisplayshortskip}{4pt plus 2pt}
\setlength{\belowdisplayshortskip}{4pt plus 2pt minus 2pt}

\theoremstyle{definition}
\newtheorem{proposition}{Proposition}

\newcommand{\R}{\mathbb R}
\newcommand{\Cont}{\mathsf C}
\newcommand{\Quit}{\mathsf Q}

\title{4PPS: A Four-Player Paired-Singleton Calibration\\
\large Two exact completions with common LCP data}
\author{}
\date{}

\begin{document}
\maketitle
\vspace{-1.5em}

\begin{abstract}
The four-player paired-singleton family (4PPS) fixes the normalized
single-quitter comparison matrix of a four-player quitting game.  This matrix
has full corrected core, is standard $Q$, has no homogeneous simplex solution,
and is not projective $\overline Q$.  Two explicit completions share it: one
has an exact pure stationary equilibrium, while the other has no stationary
exact terminal Nash equilibrium but has an exact absorbing period-two terminal
Nash profile and a uniform-equilibrium payoff.  Both lie in the same
residual-hard LCP class, and both have zero infimum of maximum positive
terminal deviation gain.
\end{abstract}

\section{The 4PPS family}

Let $I=\{0,1,2,3\}$, with paired blocks $\{0,1\}$ and $\{2,3\}$.
A quitting game assigns a payoff vector $r(S)\in\R^4$ to every nonempty
coalition $S\subseteq I$ of players who quit at the first quitting stage.
Perpetual continuation has payoff $0$.

The normalized singleton matrix of $r$ is
\begin{equation}
\label{eq:normalized-solo}
 A^r_{ij}=r_i(\{j\})-r_i(\{i\}).
\end{equation}
The 4PPS family is the class of tables for which
\begin{equation}
\label{eq:M}
 A^r=M,
 \qquad
 M=
 \begin{pmatrix}
 0&3&-1&-1\\
 3&0&-1&-1\\
 -1&-1&0&3\\
 -1&-1&3&0
 \end{pmatrix}.
\end{equation}
Thus a sole quitter exceeds the same player's own singleton benchmark by
$0$ for the quitter, by $3$ for the quitter's partner, and by $-1$ for either
member of the other pair.  Rewards for coalitions of size at least two are
not fixed by membership in 4PPS.

\section{The common matrix}

For a matrix $A$ on a finite player set, define the corrected layers by
\[
 I_0=I,
 \qquad
 I_{n+1}=\{i\in I_n:\exists j\in I_n,\ j\ne i,\ A_{ij}\le0\},
 \qquad
 I_*=\bigcap_{n\ge0}I_n.
\]
A homogeneous simplex solution is a vector $z\ge0$ satisfying
\[
 \sum_i z_i=1,
 \qquad Mz\ge0,
 \qquad z_i(Mz)_i=0\quad(i\in I).
\]
The matrix $M$ is standard $Q$ when, for every $q\in\R^4$, the
complementarity problem
\[
 w=q+Mz,
 \qquad w,z\ge0,
 \qquad z_iw_i=0\quad(i\in I)
\]
has a solution.

For a nonempty principal submatrix $M_J$, the projective problem replaces
$w=q+M_Jz$ by
\[
 w=cq+M_Jz,
 \qquad c\ge0,\quad z\ge0,\quad c+\sum_{i\in J}z_i=1,
 \qquad w\ge0,\quad z_iw_i=0.
\]
Projective $\overline Q$ means that this problem is solvable for every $q$
on every nonempty principal submatrix.

\begin{proposition}[Algebraic position of the paired matrix]
For the matrix $M$ in \eqref{eq:M},
\[
 I_*=I.
\]
Moreover, $M$ is standard $Q$, has no homogeneous simplex solution, and is
not projective $\overline Q$.
\end{proposition}

Standard $Q$ is proved by explicit complementary supports after ordering
within each pair and using pair swaps.

The last failure already occurs on $J=\{0,2\}$: for
$M_J=\bigl(\begin{smallmatrix}0&-1\\-1&0\end{smallmatrix}\bigr)$ and
$q=(-1,-1)$, the inequalities $-c-z_2\ge0$ and $-c-z_0\ge0$ force
$c=z_0=z_2=0$, contradicting $c+z_0+z_2=1$.

The residual-hard LCP class consists of tables whose corrected core is
nonempty, whose corrected-core matrix is standard $Q$ and has no homogeneous
simplex solution, and whose full normalized singleton matrix is not
projective $\overline Q$.  It is tracked as the residual half of the
full-matrix projective-$\overline Q$ split inside the necessary standard-$Q$
side; residual-hardness itself is not forced of every counterexample.

\section{An exact stationary completion}

Define $G^{(\mathrm s)}$ by the singleton rewards
\[
 r^{(\mathrm s)}(\{j\})=M_{\bullet j}
\]
and, for every coalition $S$ with $|S|\ge2$, by
\begin{equation}
\label{eq:stationary-completion}
 r^{(\mathrm s)}(S)=(-2,-2,-2,-2).
\end{equation}
Its normalized singleton matrix is $M$.

\begin{proposition}[Stationary 4PPS completion]
In $G^{(\mathrm s)}$, the pure stationary profile
\[
 (\Quit,\Cont,\Cont,\Cont)
\]
is an exact terminal Nash equilibrium against arbitrary behavioral
deviations.  Its terminal payoff is
\[
 (0,3,-1,-1),
\]
and this vector is a uniform-equilibrium payoff for the finite-horizon
average-payoff games.  The table $G^{(\mathrm s)}$ belongs to the residual-hard
LCP class.
\end{proposition}

\section{An exact period-two completion}

Define $G^{(2)}$ by the following complete terminal reward table.  The table
uses the normalization in which perpetual continuation has payoff $0$.
\begin{center}
\small
\renewcommand{\arraystretch}{1.08}
\begin{tabular}{c>{\(}l<{\)}@{\qquad}c>{\(}l<{\)}}
\toprule
$S$ & \multicolumn{1}{c}{$r^{(2)}(S)$}
    & $S$ & \multicolumn{1}{c}{$r^{(2)}(S)$}\\
\midrule
$\{0\}$ &(1,4,0,0)       &$\{1\}$ &(4,1,0,0)\\
$\{2\}$ &(0,0,1,4)       &$\{3\}$ &(0,0,4,1)\\
$\{0,1\}$ &(1,1,1,1)     &$\{0,2\}$ &(1,1,1,0)\\
$\{0,3\}$ &(1,0,1,1)     &$\{1,2\}$ &(0,1,1,1)\\
$\{1,3\}$ &(1,1,0,1)     &$\{2,3\}$ &(1,1,1,1)\\
$\{0,1,2\}$ &(1,0,0,0)   &$\{0,1,3\}$ &(0,1,0,0)\\
$\{0,2,3\}$ &(0,0,0,1)   &$\{1,2,3\}$ &(0,0,1,0)\\
$I$ &(-1,-1,-1,-1)        &&\\
\bottomrule
\end{tabular}
\end{center}
The diagonal singleton rewards all equal $1$, and subtraction of those
diagonal entries gives the same normalized matrix $M$ in \eqref{eq:M}.

A \emph{product root} is a one-stage independent mixed-action profile.  We
identify it with its vector of Quit probabilities.  For a product root $q$
and a continuation vector $v$, let $F(q;v)$ be the one-stage expected payoff,
using $r^{(2)}(S)$ when the realized quitting set $S$ is nonempty and $v$ when
every player Continues.

Let
\begin{equation}
\label{eq:parameters}
 f(x)=44x^4-7x^3-26x^2+x+3.
\end{equation}
There is a unique $a\in(1/2,1)$ with $f(a)=0$.  Select that root and set
\begin{equation}
\label{eq:secondary}
 b=\frac{4a^2-1}{3a^2}.
\end{equation}
The certified bounds are
\begin{equation}
\label{eq:parameter-bounds}
 \frac{37}{50}<a<\frac34,
 \qquad
 \frac12<b<1.
\end{equation}
Here $b>1/2$ gives $b^{-1}<2$, the phase-value bound used below.

The odd and even roots are
\begin{equation}
\label{eq:roots}
 q^{\mathrm o}=(1-a,0,1-b,0),
 \qquad
 q^{\mathrm e}=(0,1-a,0,1-b),
\end{equation}
and their continuation values are
\begin{equation}
\label{eq:values}
 v^{\mathrm o}=(1,b^{-1},1,a^{-1}),
 \qquad
 v^{\mathrm e}=(b^{-1},1,a^{-1},1).
\end{equation}
Both the displayed phase roots and the displayed continuation values are
genuinely different.
The parameter equations come directly from the odd Bellman identity
$F(q^{\mathrm o};v^{\mathrm e})=v^{\mathrm o}$: coordinate $3$ gives
$4a-3ab=a^{-1}$ and hence \eqref{eq:secondary}; coordinate $1$ gives
$b^2(3-2a)+b(1-a)-1=0$; eliminating $b$ yields
$(a-1)f(a)=0$.  The branch $a=1$ is the nonabsorbing all-Continue algebraic
branch.

\begin{proposition}[Exact two-phase Nash--Bellman block]
The displayed roots and values satisfy
\[
 F(q^{\mathrm o};v^{\mathrm e})=v^{\mathrm o},
 \qquad
 F(q^{\mathrm e};v^{\mathrm o})=v^{\mathrm e}.
\]
The Quit-minus-Continue endpoint differences at the odd and even phases are,
respectively,
\[
 \left(0, a+b-ab-b^{-1}, 0, 1-a^{-1}\right)
\]
and
\[
 \left(a+b-ab-b^{-1}, 0, 1-a^{-1}, 0\right).
\]
Both nonzero expressions are strictly negative.  The zeros occur at the
mixing players---$0,2$ in the odd phase and $1,3$ in the even phase---where
equality is required; the strict negatives occur at the sure-Continue players,
where the Nash inequality suffices.  Hence each root is exact Nash against its
displayed continuation value.  The Continue mass at each phase is $ab<1$, and
the two phases form an absorbing exact Nash--Bellman cycle.
\end{proposition}

Let $\sigma^{(2)}$ be the behavioral profile that alternates the odd and even
roots in \eqref{eq:roots}, beginning with the odd root.

\begin{proposition}[Period-two terminal and uniform equilibrium]
The profile $\sigma^{(2)}$ is an exact terminal Nash equilibrium against
arbitrary behavioral deviations, with terminal payoff
\[
 U(\sigma^{(2)})=v^{\mathrm o}=(1,b^{-1},1,a^{-1}).
\]
The same vector is a uniform-equilibrium payoff for the finite-horizon
average-payoff games.

There is no stationary product root whose stationary behavioral profile is
an exact terminal Nash equilibrium of $G^{(2)}$.  For absorbing roots this is
a finite complementarity classification; all-Continue fails because player
$0$ gains $1$ by quitting alone.  The table $G^{(2)}$ belongs to the
residual-hard LCP class.
\end{proposition}

\end{document}
